\sscp{Tree kangaroos \it{(Dendrolagus)} and wallabies}{C-08-03}
Data taken from word lists is usually unclear as to which kangaroo-like
animal the elicited term designates.
We follow the glosses in the sources,
though note that the same term is sometimes glossed by one source as ‘kangaroo/wallaby’.
Terms for ‘tree kangaroo’
and ‘kangaroo’ in SHWNG languages are given below.

\citet[222]{OsmondPawley2011} reconstruct \tsc{POc} *wakin ‘wallaby’
and point to the Moor
and Dusner words as cognates
and on this basis posit PEMP *wakin ‘wallaby’.
In addition to cognates of PEMP *wakin (or similar **waked)
two other cognate sets occur in SHWNG languages for `(tree) kangaroo': 
**musi and **amoR.

\noindent{\wg\st{2pt}\begin{xltabular}{\linewidth}
{lllll>{\raggedright\arraybackslash}X}
\lsptoprule
%Language	&	ISO	&	Term	&	Gloss	&	Etymon	&	Source	\\	\midrule
%\endfirsthead
Language	&	ISO	&	Term	&	Gloss	&	Etymon	&	Source	\\	\midrule	\hhline{~}
\endhead
As				&	asz	&	mowɔ				&	kangaroo							&	\it{misc.}&	SV \citeyear[241]{SmitsVoorhoeve1992b}	\\
Biak\footnotemark			&	bhw	&	wane 				&	(tree) kangaroo				&	*wakin		&	SV \citeyear[241]{SmitsVoorhoeve1992b}	\\

Roon			&	rnn	&	musiɛw; 		&	tree kangaroo;				&	**musi;		&	SV SV \citeyear[241]{SmitsVoorhoeve1992b};	\\	\hhline{~}
					&			&	amɔr				&	kangaroo							&	**amoR		&	SV \citeyear[263]{SmitsVoorhoeve1992b}	\\
Dusner		&	dsn	&	wa(e)n			&	kangaroo							&	*wakin		&	SV SV \citeyear[241]{SmitsVoorhoeve1992b}	\\
Meoswar		&	mvx	&	amɔr				&	kangaroo							&	**amoR		&	SV SV \citeyear[241]{SmitsVoorhoeve1992b}	\\
Yaur			&	jau	&	uahiɛ				&	kangaroo							&	\it{misc.}&	SV SV \citeyear[241]{SmitsVoorhoeve1992b}	\\
\end{xltabular}}
\footnotetext{
		\citet[263]{SmitsVoorhoeve1992b} give the following words for
		Biak under the entry for ‘tree wallaby’:  \it{wane (dame) sup}, \it{rambab} (cf.
		\it{rambav} ‘cuscus’), \it{poden sup},
		\it{bosei}, \it{wan}, \it{wansup}, \it{wantup}.
		They give the following words under their entry for ‘kangaroo’: 
		\it{wan(e) maper}, \it{wane dame mapper}, \it{wan},
		\it{keray}, \it{poden mapper}, \it{bosei},
		\it{rambab sop}.}
		
\noindent{\wg\st{2pt}\begin{xltabular}{\linewidth}
{lllll>{\raggedright\arraybackslash}X}
\lsptoprule
%Language	&	ISO	&	Term	&	Gloss	&	Etymon	&	Source	\\	\midrule
%\endfirsthead
Language	&	ISO	&	Term	&	Gloss	&	Etymon	&	Source	\\	\midrule	\hhline{~}
\endhead
Yerisiam	&	ire	&	aure;				&	tree kangaroo					&	\it{misc.}&	SV \citeyear[263]{SmitsVoorhoeve1992b}	\\	\hhline{~}
					&			&	nahiːʤaka		&	kangaroo							&	\it{misc.}&	SV \citeyear[241]{SmitsVoorhoeve1992b}	\\
Yeretuar	&	gop	&	ˈmaɛ; 			&	tree kangaroo; 				&	\it{misc.}&	SV \citeyear[263]{SmitsVoorhoeve1992b}	\\	\hhline{~}
					&			&	buka				&	kangaroo							&	\it{misc.}&	SV \citeyear[241]{SmitsVoorhoeve1992b}	\\
Moor			&	mhz	&	ˈwana(ra);	&	(tree) kangaroo;			&	*wakin		&	SV \citeyear[241]{SmitsVoorhoeve1992b}, 263	\\
					&			&	ˈwusi				&	(tree) kangaroo				&	**musi		&	SV \citeyear[241]{SmitsVoorhoeve1992b}, 263	\\
Ansus			&	and	&	(kapa) musi;&	tree kangaroo;				&	**musi		&	\cite{PriceDonohue2009}	\\
					&			&	amo					&	wallaby								&	**amoR		&	\cite{PriceDonohue2009}	\\
Ambai			&	amk	&	mansa(ŋ);		&	Grizzled Tree 				&	**mantəR;	&	\cite{Price2002}	\\	\hhline{~}
					&			&							&	Kangaroo,							&						&	\hp{\,}						\\	\hhline{~}
					&			&							&	\it{Dendrolagus}			&						&	\hp{\,}						\\	\hhline{~}
					&			&							&	\it{inustus};					&						&	\hp{\,}						\\	
					&			&	amo					&	Brown Wallaby, 				&	**amoR		&	\cite{Price2002}	\\ \hhline{~}
					&			&							&	\it{Dorcopsis veterum}&						&	\hp{\,}						\\
Wandamen	&	wad	&	musi;				&	(tree) kangaroo				&	**musi		&	SV \citeyear[241]{SmitsVoorhoeve1992b}, 263	\\
					&			&	amor				&	(tree) kangaroo				&	**amoR		&	SV \citeyear[241]{SmitsVoorhoeve1992b}, 263	\\
Serui-Laut&	seu	&	awondo;			&	tree kangaroo;				&	\it{misc.}&	SV \citeyear[263]{SmitsVoorhoeve1992b}	\\	\hhline{~}
					&			&	aˈmo				&	kangaroo							&	**amoR		&	SV \citeyear[241]{SmitsVoorhoeve1992b}	\\
Tandia		&	tni	&	dɔβiɛβiːβi	&	kangaroo							&	\it{misc.}&	SV \citeyear[241]{SmitsVoorhoeve1992b}	\\
\lspbottomrule
\end{xltabular}}
